\chapter{円柱座標系}
    \section{計量係数}
        \begin{shadebox}
            \[h_r = 1, \hph = r, h_z = 1\]
        \end{shadebox}
        \begin{proof}
            \quad\par
            位置ベクトルは
            \[\vr = \dthv{r\cosph}{r\sinph}{z}\]
            よって
            \begin{align*}
                h_r &\coloneq \norm{\partDeriv{\vr}{r}{}} = \norm{\dthv{\cosph}{\sinph}{0}} = 1\\
                \hph &\coloneq \norm{\partDeriv{\vr}{\varphi}{}} = \norm{\dthv{-r\cosph}{r\sinph}{0}} = r\\
                h_z &\coloneq \norm{\partDeriv{\vr}{z}{}} = \norm{\dthv{0}{0}{1}} = 1
            \end{align*}
        \end{proof}
    \section{単位ベクトルに関する3次元Euclid空間との関係}
        \begin{shadebox}
            \begin{align*}
                \ir &= \vix\cosph + \viy\sinph\\
                \iph &= -\vix\sinph + \viy\cosph
            \end{align*}
            \begin{align*}
                \vix &= \ir\cosph - \iph\sinph\\
                \viy &= \ir\sinph + \iph\cosph
            \end{align*}
        \end{shadebox}
        \begin{proof}
            \quad\par
            位置ベクトルは
            \[\vr = \dthv{r\cosph}{r\sinph}{z}\]
            よって
            \begin{align}
                \ir &\coloneq \frac{1}{h_r}\partDeriv{\vr}{r}{} = \dthv{\cosph}{\sinph}{0} = \vix\cosph+\viy\sinph \\
                \iph &\coloneq \frac{1}{\hph}\partDeriv{\vr}{\varphi}{} = \frac{1}{r}\dthv{-r\sinph}{r\cosph}{0} = -\vix\sinph+\viy\cosph
            \end{align}
            これを逆に解いてみる。\\
            $(1)\times\cosph-(2)\times\sinph$より$\vix = \ir\cosph - \iph\sinph$\\
            $(1)\times\sinph+(2)\times\cosph$より$\viy = \ir\sinph + \iph\cosph$\\
            を得る。
        \end{proof}
    \section{微分演算に関する直交座標系との関係}
        \begin{shadebox}
            \begin{align*}
                \partDeriv{}{x}{} &= \cosph\partDeriv{}{r}{} - \frac{\sinph}{r}\partDeriv{}{\varphi}{}\\
                \partDeriv{}{y}{} &= \sinph\partDeriv{}{r}{} + \frac{\cosph}{r}\partDeriv{}{\varphi}{}
            \end{align*}
        \end{shadebox}
        \begin{proof}
            \quad\par
            スカラー関数$\phi(x,y,z)$を微分することを考える。
            $d\phi \coloneq \phr dr + \php d\varphi + \phz dz$であるから
            \begin{align*}
                \partDeriv{\phi}{x}{} &= \phr\partDeriv{r}{x}{} + \php\partDeriv{\varphi}{x}{} + \phz\partDeriv{z}{x}{}\\
                &= \phr\ppx{\sqrt{x^2+y^2}} + \php\ppx{\varphi}
            \end{align*}
            第1項に関して
            \[\ppx{\sqrt{x^2+y^2}} = \frac{x}{\sqrt{x^2+y^2}} = \frac{r\cosph}{r} = \cosph\]
            第2項に関して
            \[\ppx{\varphi} = \ppx\atanEx{x}{y} = \frac{-y}{x^2+y^2} = \frac{-r\sin{\phi}}{r^2} = -\frac{\sin{\phi}}{r}\]
            以上より
            \[\partDeriv{}{x}{} = \cosph\partDeriv{}{r}{} - \frac{\sinph}{r}\partDeriv{}{\varphi}{}\]
            $y$に関しても全く同様にして導出できる。
        \end{proof}